\chapter{D-Dimer Test}

\section*{Definition and Overview}
The D-dimer test is a blood test that measures fragments produced when a blood clot breaks down. It is used mainly to help rule out dangerous blood clots, particularly deep vein thrombosis (DVT) in the legs and pulmonary embolism (PE) in the lungs. A normal (negative) D-dimer result makes these diagnoses very unlikely, while a raised (positive) result is not specific and can be caused by many conditions, so it triggers further investigation such as ultrasound or CT scanning. The test is quick and safe, but its value depends on being used in the right clinical context.

\section*{Epidemiology}
D-dimer testing is used very widely in emergency departments and general practice for people with possible blood clots. DVT affects around one to two people per 1,000 per year, and pulmonary embolism is a common and serious condition. Because the symptoms of clots overlap with many benign conditions, D-dimer testing helps avoid unnecessary imaging in low-risk people while ensuring that genuine clots are detected. The test is most useful when combined with a clinical assessment of the person's risk.

\section*{Aetiology and Pathophysiology}
D-dimer fragments are produced whenever the body forms and breaks down a blood clot.

\begin{itemize}
    \item \textbf{Clot formation:} the body forms a clot from fibrin when a blood vessel is injured
    \item \textbf{Clot breakdown:} an enzyme called plasmin dissolves the clot
    \item \textbf{D-dimers:} cross-linked fibrin fragments released into the blood during breakdown
    \item \textbf{Deep vein thrombosis:} a clot in a deep vein of the leg or arm
    \item \textbf{Pulmonary embolism:} a clot that travels to the lungs
    \item \textbf{Other causes of raised D-dimer:} pregnancy, surgery, infection, inflammation, cancer, trauma, and recent bleeding
    \item \textbf{Age:} D-dimer levels rise naturally with age
\end{itemize}

\section*{Clinical Features}
The test is performed when there are features that could indicate a blood clot.

\begin{itemize}
    \item \textbf{DVT symptoms:} swelling, pain, warmth, or redness of one leg or calf
    \item \textbf{PE symptoms:} sudden shortness of breath, chest pain, coughing up blood, or collapse
    \item \textbf{Risk factors:} recent surgery, immobility, cancer, pregnancy, the contraceptive pill, and previous clots
    \item \textbf{The blood test:} a simple blood draw from the arm
    \item \textbf{Results:} usually available within an hour or two
\end{itemize}

\section*{Differential Diagnosis}
The symptoms that prompt D-dimer testing have many causes.

\begin{itemize}
    \item Leg swelling from muscle injury, varicose veins, or heart failure
    \item Shortness of breath from pneumonia, asthma, or heart disease
    \item Chest pain from musculoskeletal causes, reflux, or pericarditis
    \item Blood clots versus superficial thrombophlebitis
    \item Pregnancy-related leg swelling
    \item Other causes of raised D-dimer, such as infection, inflammation, and surgery
\end{itemize}

\section*{When to Seek Medical Care}
See a doctor promptly if you have symptoms that could indicate a blood clot, especially sudden breathlessness, chest pain, or a swollen painful leg.

\textbf{Call 000 (or local emergency services) for:}
\begin{itemize}
    \item Sudden difficulty breathing or severe chest pain
    \item Coughing up blood
    \item Collapse or fainting
    \item A very swollen, painful, hot leg with sudden onset
    \item Signs of a stroke: sudden weakness, numbness, or difficulty speaking
\end{itemize}

\section*{Diagnosis and Investigations}
D-dimer is used as part of a structured diagnostic pathway.

\begin{itemize}
    \item \textbf{Clinical risk assessment:} scores such as the Wells score estimate the pre-test probability of a clot
    \item \textbf{D-dimer test:} used when the clinical risk is low or moderate
    \item \textbf{Low D-dimer:} makes a clot very unlikely, and no further testing is needed
    \item \textbf{High D-dimer:} leads to confirmatory imaging
    \item \textbf{Leg ultrasound:} compression ultrasound for suspected DVT
    \item \textbf{CT pulmonary angiogram:} CT scan of the lung arteries for suspected PE
    \item \textbf{High-risk presentations:} go straight to imaging, as a normal D-dimer cannot be trusted in high-risk situations
\end{itemize}

\section*{Management}
The D-dimer test itself needs no treatment; it guides further care.

\begin{itemize}
    \item \textbf{Negative D-dimer with low risk:} a clot is effectively excluded, and no anticoagulation is needed
    \item \textbf{Positive D-dimer with confirmed clot:} anticoagulant treatment with heparin, then warfarin or a DOAC
    \item \textbf{Positive D-dimer without a clot:} the cause of the raised level is considered, and treatment targets the underlying condition
    \item \textbf{Positive D-dimer with negative imaging:} the person is reviewed and other causes investigated
    \item \textbf{Follow-up:} repeat assessment if symptoms persist or recur
\end{itemize}

\section*{Complications}
\begin{itemize}
    \item False-positive results leading to unnecessary imaging and anxiety
    \item False-negative results if the test is used inappropriately in high-risk people
    \item Delayed diagnosis if a clot is missed
    \item Complications of unnecessary anticoagulation
    \item The need for repeat tests when results are borderline
    \item Radiation exposure from CT scanning when the test is positive
\end{itemize}

\section*{Prognosis and Outcomes}
When used correctly, the D-dimer test is highly effective at ruling out blood clots, allowing many people to avoid unnecessary imaging and anticoagulation. Its sensitivity is high, meaning few genuine clots are missed, and its specificity is limited, meaning positive results are common and require imaging. With the right clinical assessment, D-dimer testing improves outcomes by detecting clots early while avoiding harm from over-investigation. Treated DVT and PE carry a good prognosis, with a low risk of recurrence after the initial treatment period.

\section*{Prevention and Self-Care}
\begin{itemize}
    \item Seek prompt medical assessment for possible clot symptoms
    \item Move regularly during long travel and after surgery, and stay hydrated
    \item Follow compression and anticoagulant treatment as prescribed
    \item Tell the doctor about risk factors, including the contraceptive pill and pregnancy
    \item Attend follow-up imaging and appointments
    \item Understand that a raised D-dimer has many causes and does not by itself mean you have a clot
    \item Know the signs of a clot and act promptly
\end{itemize}

\section*{Sources and Further Reading}
The following reputable sources provide further information for healthcare students and patients seeking reliable, current advice about the D-dimer test and blood clots.

\begin{itemize}
    \item Healthdirect Australia. \textit{D-dimer test.} https://www.healthdirect.gov.au/
    \item Royal College of Pathologists of Australasia. \textit{D-dimer pathology information.} https://www.rcpa.edu.au/
    \item Lab Tests Online AU. \textit{D-dimer.} https://labtestsonline.org.au/
    \item MedlinePlus. \textit{D-dimer test.} https://medlineplus.gov/ency/article/003582.htm
    \item Australian Government Department of Health and Aged Care. \textit{Venous thromboembolism resources.} https://www.health.gov.au/
    \item MSD Manual. \textit{Diagnosis of deep vein thrombosis.} https://www.msdmanuals.com/
\end{itemize}
